% ============================================================
%  SECTION — BUDGET
%
%  HOW TO CUSTOMISE
%  ─────────────────
%  1. Fill in each line item: description, quantity, unit cost, total.
%  2. The "Total" row is computed manually — update it when you add rows.
%  3. Add/remove \midrule-separated categories as needed.
%  4. Change "MYR" to your currency if necessary.
%  5. Note funding source / acquisition method in the footnote line.
% ============================================================



\begin{frame}{Budget}
  \centering
  \small

  % X column stretches to fill \textwidth, keeping numeric columns fixed-width
  \begin{tabularx}{\textwidth}{X r r r}
    \toprule
    \textbf{Item} & \textbf{Qty} & \textbf{Unit (MYR)} & \textbf{Total (MYR)} \\
    \midrule
    % ---- Equipment ----
    \multicolumn{4}{l}{\textit{Equipment}} \\
    TODO: Item name  & TODO & TODO & TODO \\
    TODO: Item name  & TODO & TODO & TODO \\
    \midrule
    % ---- Materials / Consumables ----
    \multicolumn{4}{l}{\textit{Materials \& Consumables}} \\
    TODO: Item name  & TODO & TODO & TODO \\
    TODO: Item name  & TODO & TODO & TODO \\
    \midrule
    % ---- Software / Licences ----
    \multicolumn{4}{l}{\textit{Software \& Licences}} \\
    Open-source only & -- & -- & 0 \\
    \midrule
    % ---- Miscellaneous ----
    \multicolumn{4}{l}{\textit{Miscellaneous}} \\
    TODO: Item name  & TODO & TODO & TODO \\
    \midrule[\heavyrulewidth]
    \textbf{Total}   &      &      & \textbf{TODO} \\
    \bottomrule
  \end{tabularx}

  \vspace{0.5em}
  % Funding source, acquisition notes (purchase / lab stock / borrowed), spending cap
  {\footnotesize TODO: e.g.\ Requested from Physics Dept.\ student project fund (cap: MYR~500).}
\end{frame}